%! Choosing and Comparing Models in Context — Unit 2, Lesson 2.7 — Slides (Beamer)
\documentclass[aspectratio=169, 11pt]{beamer}
\usepackage{atda-beamer}   % provides \royalheader{} and \sectionlabel[color]{}
\usepackage{pgfplots}      % atda-beamer does NOT load pgfplots
\pgfplotsset{compat=1.18}

\begin{document}

% TITLE SLIDE (CourseName is not defined in beamer, so the name is literal)
{
\setbeamercolor{background canvas}{bg=royal}
\begin{frame}[plain]
  \vspace{2.8cm}\hspace{0.55cm}%
  \begin{minipage}{0.9\textwidth}
    {\color{goldacc}\bfseries\small LESSON 2.7}\\[0.5cm]
    {\color{white}\bfseries\LARGE Choosing and Comparing Models in Context}\\[1.2cm]
    {\color{mistmid}\small Algebra, Trigonometry, and Data Analysis~$\cdot$~Unit 2}
  \end{minipage}
\end{frame}
}

% ── HOOK ──────────────────────────────────────────────────────────────────────
\begin{frame}[t, plain]
  \royalheader{Two Students, One Table, and Both Are Right About What They Checked}
  \sectionlabel[goldacc]{hook --- a school video, views in thousands}
  \begin{center}
  \renewcommand{\arraystretch}{1.25}
  \begin{tabular}{@{}|l|c|c|c|c|@{}}
  \hline
  Day $d$ & $0$ & $1$ & $2$ & $3$ \\ \hline
  Views $V$ (thousands) & $3$ & $6$ & $12$ & $24$ \\
  \hline
  \end{tabular}
  \end{center}
  \begin{block}{Everything both of them said about these numbers is true}
    \small
    \textbf{A:} ``The jumps are $3$, then $6$, then $12$. They keep getting bigger, and a parabola's
    jumps get bigger too. \textbf{So it's quadratic.}''\\[0.15cm]
    \textbf{B:} ``Every day is double the day before. \textbf{So it's exponential.}''
  \end{block}
  \vspace{0.1cm}
  \textbf{So one of them checked something that was not \emph{enough}. Which one?}
\end{frame}

% ── THE RESOLUTION ────────────────────────────────────────────────────────────
\begin{frame}[t, plain]
  \royalheader{``The Jumps Get Bigger'' Rules Out Linear --- and Nothing Else}
  \sectionlabel{that sentence is the whole lesson}
  \begin{block}{Student B answered it}
    \centering\large
    \textbf{Second differences $3, 6$: not constant. \quad Ratios $2, 2, 2$: constant.}
  \end{block}
  \begin{itemize}
    \setlength{\itemsep}{0.3cm}
    \item \emph{Every} family except linear has growing first differences --- so ``it speeds up''
          cannot tell quadratic from exponential.
    \item You have to check which row comes out \textbf{constant}, not which row is growing.
    \item \textbf{Never name a family without naming the test.}
  \end{itemize}
\end{frame}

% ── THE TWO TESTS ─────────────────────────────────────────────────────────────
\begin{frame}[t, plain]
  \royalheader{Two Tests, Run in Order}
  \sectionlabel{\texttt{AFDA.AF.1c} --- run them on a table with \emph{equal} input steps}
  \renewcommand{\arraystretch}{1.35}
  \small
  \begin{center}
  \begin{tabular}{@{}|l|c|c|c|c||l|@{}}
  \hline
  \textbf{(a) rideshare fare} & $3$ & $5$ & $7$ & $9$ & first differences $+2, +2, +2$ \\ \hline
  \textbf{(b) braking distance} & $20$ & $45$ & $80$ & $125$ & second differences $+10, +10$ \\ \hline
  \textbf{(c) a shared post} & $5$ & $15$ & $45$ & $135$ & ratios $3, 3, 3$ \\
  \hline
  \end{tabular}
  \end{center}
  \vspace{0.3cm}
  \begin{enumerate}
    \setlength{\itemsep}{0.2cm}
    \item First differences constant $\Rightarrow$ \textbf{linear.} \quad (a)
    \item Not those, but second differences constant $\Rightarrow$ \textbf{quadratic.} \quad (b)
    \item Ratios constant $\Rightarrow$ \textbf{exponential.} \quad (c)
  \end{enumerate}
  \vspace{0.15cm}
  \textbf{The condition, and it is not optional: the inputs must go up in equal steps.} (b) steps by
  $10$, and that is fine. A table stepping $1, 2, 5, 10$ is not.
\end{frame}

% ── THE COMPARISON SHEET ──────────────────────────────────────────────────────
\begin{frame}[t, plain]
  \royalheader{The Same Three Families, Five Ways}
  \sectionlabel{\texttt{AFDA.AF.1g} and \texttt{AF.2h} --- the summary sheet for the whole unit}
  \renewcommand{\arraystretch}{1.3}
  \small
  \begin{tabularx}{\textwidth}{@{}p{2.6cm}|X|X|X@{}}
    & \textbf{Linear} & \textbf{Quadratic} & \textbf{Exponential} \\
    \hline
    Parent (2.2) & $f(x)=x$ & $g(x)=x^2$ & $h(x)=b^{\,x}$ \\
    Table test & first differences & second differences & ratios \\
    Graph shape & straight edge & one turning point & no turn; it levels off \\
    Max, min (2.4) & neither & one, at the vertex & neither \\
    Asymptote (2.5) & none & none & exactly one, horizontal \\
    The story says & same \emph{amount} added & it turns around & same \emph{factor} multiplied \\
  \end{tabularx}
  \vspace{0.3cm}

  \textbf{``Plus \$30 a month'' and ``up $30\%$ a month'' sound almost the same out loud --- and they
  are two different families.} Look for a turning point first; if it turns, you are done.
\end{frame}

% ── USING THE MODEL ───────────────────────────────────────────────────────────
\begin{frame}[t, plain]
  \royalheader{Using the Model: Algebraically \emph{and} Graphically}
  \sectionlabel{\texttt{AFDA.AF.1e} --- the video model is $V(d) = 3 \cdot 2^{\,d}$}
  \begin{center}
  \begin{tikzpicture}
  \begin{axis}[
      width=9.2cm, height=3.4cm,
      xlabel={$d$ (days)}, ylabel={$V$ (thousands)},
      xmin=0, xmax=6, ymin=0, ymax=200,
      xtick={0,1,2,3,4,5,6}, ytick={0,50,100,150,200}, minor x tick num=1,
      grid=both, grid style={linegray!45}, minor grid style={linegray!30},
      axis lines=left, tick label style={font=\tiny}, label style={font=\tiny}]
    \addplot[goldacc, thick, dashed] coordinates {(0,50) (6,50)};
    \addplot[royal, very thick, domain=0:6, samples=80]{3*2^x};
  \end{axis}
  \end{tikzpicture}
  \end{center}
  \vspace{-0.35cm}
  \small
  \textbf{Algebraically:} $V(5) = 3 \cdot 2^5 = 3 \cdot 32 = 96$ thousand views.\\[0.1cm]
  \textbf{Graphically:} across from $50$, down to the axis --- it passes $50$ thousand between day
  $4$ and day $5$, because $V(4)=48$ and $V(5)=96$.\\[0.1cm]
  \textbf{Name which one you are doing.} The standard asks for both.
\end{frame}

% ── THE WINDOW ────────────────────────────────────────────────────────────────
\begin{frame}[t, plain]
  \royalheader{$3.1$ Billion Views --- Is the Model Wrong?}
  \sectionlabel[goldacc]{the same model, read at day $20$}
  \begin{block}{$V(20) = 3 \cdot 2^{20} = 3{,}145{,}728$ thousand $\approx 3.1$ billion views}
    \centering\large
    \textbf{No. The model is being used outside the days the data covers.}
  \end{block}
  \begin{itemize}
    \setlength{\itemsep}{0.28cm}
    \item Predicting beyond the inputs your data covered is called \textbf{extrapolation}.
    \item Doubling has to stop --- the school runs out of \emph{people}.
    \item \textbf{A model is a choice, not a fact, and every model has a window where it works.}
  \end{itemize}
  \vspace{0.15cm}
  Say what the model predicts \emph{and} whether you believe it.
\end{frame}

% ── GUIDED PRACTICE ───────────────────────────────────────────────────────────
\begin{frame}[t, plain]
  \royalheader{Guided Practice: Three Measurements, Three Families}
  \sectionlabel{a draining pool, a punted football, a dose of caffeine}
  \renewcommand{\arraystretch}{1.3}
  \small
  \begin{center}
  \begin{tabular}{@{}|l|c|c|c|c|c||l|@{}}
  \hline
  (a) water left (gal) & $4000$ & $3650$ & $3300$ & $2950$ & $2600$ & \textbf{linear} \\ \hline
  (b) ball height (ft) & $0$ & $48$ & $64$ & $48$ & $0$ & \textbf{quadratic} \\ \hline
  (c) caffeine (mg) & $160$ & $80$ & $40$ & $20$ & $10$ & \textbf{exponential} \\
  \hline
  \end{tabular}
  \end{center}
  \vspace{0.3cm}
  \textbf{(b)} maximum $64$ feet at $t=2$; zeros at $t=0$ and $t=4$ --- the punt and the landing.\\[0.1cm]
  \textbf{(c)} $C(25) = 160 \cdot (0.5)^5 = 5$ mg, and it \emph{never} reaches zero --- horizontal
  asymptote, Lesson 2.5.\\[0.1cm]
  \textbf{Item 4 is the one that earns the time:} (a) and (c) both go down. With only two readings of
  each, \textbf{you cannot tell them apart.}
\end{frame}

% ── TWO POINTS ARE NEVER ENOUGH ───────────────────────────────────────────────
\begin{frame}[t, plain]
  \royalheader{Two Points Are Never Enough}
  \sectionlabel{$4$ cells on day $0$, $8$ cells on day $1$ --- both models fit exactly}
  \begin{center}
  \begin{tikzpicture}
  \begin{axis}[
      width=8.6cm, height=3.4cm,
      xlabel={$x$ (days)}, ylabel={cells},
      xmin=0, xmax=3, ymin=0, ymax=36,
      xtick={0,1,2,3}, ytick={0,8,16,24,32}, minor x tick num=1,
      grid=both, grid style={linegray!45}, minor grid style={linegray!30},
      axis lines=left, tick label style={font=\tiny}, label style={font=\tiny}]
    \addplot[royal, very thick] coordinates {(0,4) (3,16)};
    \addplot[goldacc, very thick, domain=0:3, samples=60]{4*2^x};
    \addplot[royal, only marks, mark=*, mark size=2pt] coordinates {(0,4) (1,8)};
  \end{axis}
  \end{tikzpicture}
  \end{center}
  \vspace{-0.35cm}
  \small
  \textbf{Blue: $4+4x$. \quad Gold: $4 \cdot 2^{x}$.} They agree perfectly on every day the biologist
  measured.\\[0.1cm]
  \textbf{The question is not ``which model?'' --- it is ``what would you go and measure?''} Count the
  cells on day $2$: linear says $12$, exponential says $16$.
\end{frame}

% ── ACTIVITY LAUNCH ───────────────────────────────────────────────────────────
\begin{frame}[t, plain]
  \royalheader{Group Activity: Which Family Owns This Situation?}
  \sectionlabel{groups of three --- everyone starts at Tier R}
  \begin{block}{The rule for the whole activity}
    \textbf{Never name a family without naming the test.} From a table that is a \emph{row}; from a
    graph it is a \emph{feature} --- a straight edge, a turning point, a flattening end.
  \end{block}
  \vspace{0.1cm}
  \begin{itemize}
    \setlength{\itemsep}{0.25cm}
    \item \textbf{Tier R:} two patches of duckweed --- one adds $4$ m$^2$ a week, one triples. Find
          the week the multiplier passes the adder.
    \item \textbf{Tier A:} a \$17 fare read two ways, and a football model that says $-960$ feet.
    \item \textbf{Tier E:} a coffee cup whose ratios are \emph{not} constant --- until you subtract
          the room.
  \end{itemize}
\end{frame}

% ── CLOSE ─────────────────────────────────────────────────────────────────────
\begin{frame}[t, plain]
  \royalheader{Exit Ticket}
  \sectionlabel[goldacc]{notes closed --- 5 minutes}
  A gym counted its members: $64$, $96$, $144$, $216$, $324$ over five months.
  \begin{enumerate}
    \setlength{\itemsep}{0.3cm}
    \item Fill the three test rows. Name the family --- \textbf{and name the row that decided it}.
    \item Predict months $5$ and $6$. The building holds $500$. How far ahead can this model be
          trusted?
    \item Somebody wrote ``the gains speed up, so it's quadratic.'' Noticing the growing jumps was
          \emph{correct} in item 1. Say what is wrong anyway.
  \end{enumerate}
  \vspace{0.2cm}
  {\color{royal}\bfseries That is Unit 2. Next: the unit test, then Unit 3 --- Exponential and
  Logarithmic Functions.}
\end{frame}

\end{document}
